\section{Related Work}

\subsection{Dynamic and Traffic-Aware Routing}

Route choice under congestion rests on the equilibrium tradition founded by
Wardrop~\cite{wardrop1952}. Reactive guidance --- routing on recently measured
link travel times --- is the operational descendant of that tradition, and its
central difficulty is well documented: information supplied to drivers need not
reduce congestion, because the response to the information changes the state
it described~\cite{arnott1991information}. Acemoglu et al.~\cite{acemoglu2018braess}
give this a sharp game-theoretic form, showing that additional information can
strictly worsen equilibrium outcomes. The severity depends on how much of the
demand responds, which makes guidance penetration a first-order variable rather
than an implementation detail~\cite{yang1998penetration}.

Two policy families respond to these failures directly. Load-balancing guidance
distributes vehicles over a set of travel-time-acceptable alternatives according
to predicted link load; we use the flow-balanced $k$-shortest-path construction
of Pan et al.~\cite{pan2013proactive}. Reliability-aware guidance adds a
dispersion term to expected travel time, following the mean-variance
formulation of Sen et al.~\cite{sen2001meanvariance}. Both are established
families with published objectives, and we take them as given.

\subsection{Environmental, Multi-Criteria, and Context-Dependent Routing}

Eco-routing minimises an emissions proxy rather than
time~\cite{boriboonsomsin2012eco}, and raises the question of whether an
environmental objective is \emph{distinct} from a temporal one in a given
network. That is an empirical property of the environment, not a property of
the objective, and we treat it as a measurement rather than an assumption
(Section~\ref{sec:pareto}).

More generally, comparing routing policies on a single scalar hides the
structure that multiple criteria expose. Where several criteria disagree, the
identity of the preferred policy depends on which is chosen, and a Pareto
analysis makes that dependence visible before any weighting is applied. We
adopt that ordering --- dominance first, scalarisation only as sensitivity ---
and avoid monetary valuation entirely, since no verified value of time, value
of reliability or carbon price was available for this setting.

\subsection{Algorithm Selection and Decision-Oriented Evaluation}

Choosing among algorithms by instance features is Rice's algorithm-selection
problem~\cite{rice1976}, with a mature methodology developed largely in
combinatorial search~\cite{xu2008satzilla,smithmiles2009,kotthoff2014survey}
and standard reference points --- the virtual best and single best solver ---
from benchmarking practice~\cite{bischl2016aslib}. Instance-space analysis asks
where in a feature space each algorithm is strong~\cite{smithmiles2014instance},
which is precisely the form RQ1 takes for routing policies.

Two further strands matter here. First, evaluating a predictor by the cost of
the decision it induces rather than by its predictive error is the
predict-then-optimise distinction of Elmachtoub and
Grigas~\cite{elmachtoub2022spo}; our results show the two criteria ranking
methods differently. Second, declining to decide is
classical~\cite{chow1970reject}, and split conformal prediction supplies
distribution-free intervals to decide when~\cite{vovk2005alrw,lei2018conformal}.
Its coverage guarantee rests on exchangeability, which covariate shift
breaks~\cite{shimodaira2000,tibshirani2019covshift} --- a limitation we treat
as a design constraint rather than a caveat.

\subsection{Positioning}
\label{sec:positioning}

The literature supplies the routing objectives, the traffic-feedback mechanisms
that make them fail, the algorithm-selection framework, the regret and
virtual-best reference points, and the uncertainty tools --- each separately and
each well developed. This study combines them into a single empirical
policy-selection evaluation, run under a controlled factorial of traffic
conditions in a physical microsimulation, and reports decision regret rather
than selection accuracy as the primary quantity.

We do not claim that regime-dependent routing has not been studied, nor that no
prior work has combined any subset of these ideas. What we offer is the
measurement: how much decision value the combination actually carries in a
setting where every component is held fixed and known.
